% 3-3-3-eda.tex
%  Budget: 3pp.  EDA -- figures 01-09, already exported to reports/figures/.
%  Source map: multiple seasonality + volatility clustering + heteroscedasticity
%   as intrinsic features of electricity price series -> jiang_review_2018 ;
%   renewables penetration driving negative prices and added volatility
%   -> oconnor_review_2025 ; lag structure motivating the D-1/D-2/D-3/D-7
%   feature convention -> lago_forecasting_2021 .
%  All statistics verified 2026-09-05 against data/raw/DE.csv:
%   hourly profile min h3=22.78 max h19=46.01; weekday Sun 23.325 Sat 30.075
%   Tue 38.701; yearly 2012:42.890 ... 2016:28.979 2017:34.201;
%   autocorr lag1 0.9337 lag24 0.6598 lag168 0.5855;
%   corr(price,exog_1)=+0.6109  corr(price,exog_2)=-0.3885;
%   daily baseload n=2184 mean 34.6779 std 11.9713 min -56.87 max 101.9217.
%  NOTE: the 3-sigma value is NOT quoted anywhere here -- retired as a regime
%   switch (HANDOFF 19.0.1); k=1.5 -> 62.6989 is introduced in 3-8, not here.

\subsection{تحلیل اکتشافی داده}
\label{sec:3-3-3}

هدف این بخش توصیف داده به‌خاطر خودِ توصیف نیست. هر یک از شش الگوی زیر، تصمیمی
روش‌شناختی را در بخش‌های بعدی توجیه می‌کند، و همین پیوند است که تحلیل اکتشافی را از
یک ضمیمه‌ی تشریفاتی جدا می‌سازد.

\subsubsection*{توزیع قیمت}

\begin{figure}[htbp]
  \centering
  \includegraphics[width=0.85\textwidth]{01_price_distribution}
  \caption{توزیع قیمت‌های ساعتی بازار روز-پیش \lr{EPEX-DE} در دوره‌ی ۲۰۱۲ تا ۲۰۱۷،
  شامل ناحیه‌ی قیمت‌های منفی.}
  \label{fig:price-dist}
\end{figure}

\cref{fig:price-dist} توزیع قیمت‌های ساعتی را نشان می‌دهد. توده‌ی اصلی توزیع در بازه‌ی تقریبی ۲۰ تا
۵۰ یورو بر مگاوات‌ساعت متمرکز است، اما دنباله‌ی چپ تا ناحیه‌ی منفی امتداد می‌یابد و
دنباله‌ی راست تا بیش از دو برابر میانگین. کشیدگی \lr{16.27} در برابر مقدار سه برای توزیع
نرمال، سنجه‌ای عددی از همین دنباله‌های سنگین است.

قیمت‌های منفی \lr{—} ۵۳۸ ساعت، معادل \lr{1.03} درصد کل دوره \lr{—} پیامد ساختاری نفوذ
بالای تجدیدپذیرها هستند، نه نوفه‌ی داده \cite{oconnor_review_2025}. این ویژگی دو
پیامد روش‌شناختی دارد که در \cref{sec:3-5} به آن‌ها بازمی‌گردیم: کنارگذاشتن \lr{MAPE}، و
انتخاب معیارهایی که به علامت مخرج حساس نباشند.

\subsubsection*{فصلی‌بودن چندگانه}

\begin{figure}[htbp]
  \centering
  \includegraphics[width=0.80\textwidth]{02_hourly_seasonality}
  \caption{نیم‌رخ درون‌روزیِ قیمت: میانگین قیمت به تفکیک ساعت شبانه‌روز.}
  \label{fig:hourly-seasonality}
\end{figure}

\begin{figure}[htbp]
  \centering
  \includegraphics[width=0.80\textwidth]{03_weekly_seasonality}
  \caption{نیم‌رخ هفتگیِ قیمت: میانگین قیمت به تفکیک روز هفته.}
  \label{fig:weekly-seasonality}
\end{figure}

\begin{figure}[htbp]
  \centering
  \includegraphics[width=0.80\textwidth]{04_annual_seasonality}
  \caption{نیم‌رخ سالانه‌ی قیمت در سراسر دوره‌ی مطالعه.}
  \label{fig:annual-seasonality}
\end{figure}

سری قیمت برق هم‌زمان در سه مقیاس فصلی است، و
\crefrange{fig:hourly-seasonality}{fig:annual-seasonality} هر سه را نشان
می‌دهند. این چندلایگی از خصوصیات شناخته‌شده‌ی این سری است \cite{jiang_review_2018}.

در مقیاس شبانه‌روزی، میانگین قیمت از کمینه‌ی \lr{22.78} در ساعت سه بامداد تا بیشینه‌ی \lr{46.01}
در ساعت نوزده نوسان می‌کند \lr{—} یعنی بیش از دو برابر. الگو دو قلّه دارد: یک قلّه‌ی
صبحگاهی حوالی ساعت هشت (\lr{42.20}) و قلّه‌ی اصلیِ عصرگاهی. همین دامنه‌ی وسیع، دلیل تصمیمی
است که در بخش‌های ۳-۶ و ۳-۷ تکرار می‌شود: به‌جای یک مدل واحد برای همه‌ی ساعت‌ها،
به‌ازای هر ساعتِ شبانه‌روز یک مدل مستقل برازش می‌شود، چون رفتار آماری ساعت سه بامداد
و ساعت نوزده یکی نیست.

در مقیاس هفتگی، تفاوت میان روزهای کاری و تعطیل چشمگیر است: میانگین یکشنبه \lr{23.33} و
شنبه \lr{30.08} در برابر \lr{38.70} برای سه‌شنبه، پرقیمت‌ترین روز هفته. این شکاف، هم متغیرهای
مجازیِ روز هفته در ماتریس ویژگی (\cref{sec:3-4}) و هم قاعده‌ی رفتاریِ مدل ساده‌ی مرجع را
توجیه می‌کند؛ آن قاعده برای دوشنبه به آخرین روز کاری و برای تعطیلات به هم‌تای هفته‌ی
پیشین ارجاع می‌دهد، دقیقاً چون میانگین یکشنبه از میانگین دوشنبه بسیار دور است.

در مقیاس سالانه، روند نزولیِ ۲۰۱۲ تا ۲۰۱۶ و بازگشت سال ۲۰۱۷ که در \cref{sec:3-3-1} گزارش
شد، در \cref{fig:annual-seasonality} دیده می‌شود.

\subsubsection*{خوشه‌بندی نوسان و مقادیر حدی}

\begin{figure}[htbp]
  \centering
  \includegraphics[width=0.85\textwidth]{05_volatility_clustering}
  \caption{خوشه‌بندی نوسان: دوره‌های آرام و پرتلاطم به‌صورت خوشه‌ای در زمان
  کنار هم می‌نشینند.}
  \label{fig:vol-clustering}
\end{figure}

\cref{fig:vol-clustering} خوشه‌بندی نوسان را نشان می‌دهد: دوره‌های آرام و دوره‌های پرتلاطم به‌صورت
خوشه‌ای کنار هم می‌نشینند و به‌طور تصادفی در زمان پراکنده نیستند. این ناهمگنیِ
واریانس، ویژگیِ شناخته‌شده‌ی سری قیمت برق است \cite{jiang_review_2018}.

مسیر از همین مشاهده مستقیماً به \cref{sec:3-8} می‌رسد. اگر رفتار بازار میان دوره‌های آرام و
پرتنش به‌قدر کافی متفاوت باشد، مجموعه‌ی وزنی که در حالت آرام بهینه است لزوماً در حالت
پرتنش بهینه نیست \lr{—} و همین، منطقِ وزن‌دهیِ گروهیِ رژیم‌محور است. تعریف عملیاتیِ
دو رژیم و آستانه‌ی تفکیک آن‌ها در \cref{sec:3-8} می‌آید، نه اینجا: آستانه بر پایه‌ی داده‌ی
اعتبارسنجی و با قاعده‌ای تصریح‌شده تعیین می‌شود، و استخراج آن از یک نمودار اکتشافی
همان انتخابِ درون‌نمونه‌ای است که پژوهش از آن پرهیز می‌کند.

\subsubsection*{ساختار خودهمبستگی}

\begin{figure}[htbp]
  \centering
  \includegraphics[width=0.85\textwidth]{06_acf_pacf}
  \caption{توابع خودهمبستگی و خودهمبستگی جزئی سری قیمت ساعتی.}
  \label{fig:acf-pacf}
\end{figure}

\cref{fig:acf-pacf} توابع خودهمبستگی
(\lr{ACF}\footnote{\lr{Autocorrelation Function}}) و خودهمبستگی جزئی
(\lr{PACF}\footnote{\lr{Partial Autocorrelation Function}}) را نشان می‌دهد. خودهمبستگی در وقفه‌ی یک
ساعت \lr{0.9337}، در وقفه‌ی ۲۴ ساعت \lr{0.6598} و در وقفه‌ی ۱۶۸ ساعت (یک هفته) \lr{0.5855} است.

قلّه‌های آشکار در وقفه‌های ۲۴ و ۱۶۸ ساعت، دو دوره‌ی تناوبِ روزانه و هفتگی را تأیید
می‌کنند و مبنای دو تصمیم‌اند. نخست، انتخاب وقفه‌های قیمتی ماتریس ویژگی \lr{—} روزهای
\lr{D-1}، \lr{D-2}، \lr{D-3} و \lr{D-7} \lr{—} که همان قرارداد مرجع است
\cite{lago_forecasting_2021} و اینجا شاهد تجربیِ خود را در داده می‌یابد. دوم، انتخاب
مرتبه‌ی فصلیِ هفت‌روزه برای مدل \lr{SARIMAX} در \cref{sec:3-6}: بُعد شبانه‌روزی پیشاپیش با
تفکیک ۲۴گانه‌ی مدل‌ها جذب شده است، پس آنچه برای مؤلفه‌ی فصلی باقی می‌ماند تناوب هفتگی
است.

\subsubsection*{متغیرهای برون‌زا}

\begin{figure}[htbp]
  \centering
  \includegraphics[width=0.85\textwidth]{07_exogenous_correlation}
  \caption{همبستگی قیمت با دو متغیر برون‌زا: پیش‌بینی روزپیشِ بار و پیش‌بینی
  روزپیشِ تولید بادی و خورشیدی.}
  \label{fig:exog-corr}
\end{figure}

\begin{figure}[htbp]
  \centering
  \includegraphics[width=0.85\textwidth]{07b_rolling_correlation}
  \caption{همبستگی غلتانِ قیمت با متغیرهای برون‌زا در طول زمان؛ شدت رابطه ثابت
  نمی‌ماند.}
  \label{fig:rolling-corr}
\end{figure}

\cref{fig:exog-corr} رابطه‌ی قیمت با دو متغیر برون‌زا را نشان می‌دهد و علامت‌ها همان است که نظریه‌ی
دستور تولید انتظار دارد. همبستگی قیمت با پیش‌بینی بار مثبت و برابر \lr{0.6109} است:
تقاضای بیشتر، واحدهای گران‌تر را وارد مدار می‌کند. همبستگی آن با پیش‌بینی تولید بادی
و خورشیدی منفی و برابر \lr{0.3885} است: تولید تجدیدپذیر با هزینه‌ی نهاییِ نزدیک به صفر،
منحنی عرضه را به سمت راست می‌راند و قیمت تسویه را پایین می‌آورد.

\cref{fig:rolling-corr} همین همبستگی‌ها را به‌صورت غلتان در طول زمان نشان می‌دهد و روشن می‌کند که
شدت آن‌ها ثابت نیست. این ناپایداریِ رابطه، یکی از دلایلی است که واسنجی مجدد روزانه را
ضروری می‌کند: ضرایبی که بر داده‌ی سال ۲۰۱۳ برازش شده‌اند، برای سال ۲۰۱۷ همان ضرایب
درست نیستند.

\subsubsection*{شکست‌های ساختاری و سری روزانه}

\begin{figure}[htbp]
  \centering
  \includegraphics[width=0.90\textwidth]{08_full_series_structural_breaks}
  \caption{کل سری قیمت ساعتی به‌همراه شکست‌های ساختاریِ شناسایی‌شده.}
  \label{fig:structural-breaks}
\end{figure}

\begin{figure}[htbp]
  \centering
  \includegraphics[width=0.90\textwidth]{09_daily_baseload}
  \caption{سری بار پایه‌ی روزانه: میانگین ۲۴ قیمت ساعتی هر روز.}
  \label{fig:daily-baseload}
\end{figure}

\cref{fig:structural-breaks} کل سری را همراه با شکست‌های ساختاریِ آن نشان می‌دهد. وجود این شکست‌ها،
تصمیمِ استفاده از پنجره‌ی واسنجیِ غلتانِ ۱۰۹۲روزه به‌جای آموزش بر کل تاریخچه را توجیه
می‌کند: پنجره‌ی غلتان، داده‌ی کهنه‌ی متعلق به رژیم‌های پیشین بازار را به‌تدریج کنار
می‌گذارد.

\cref{fig:daily-baseload} سری بار پایه‌ی روزانه \lr{—} یعنی میانگین ۲۴ قیمت ساعتی هر روز \lr{—} را نشان
می‌دهد. این سری \lr{2{,}184} روز دارد، با میانگین \lr{34.68}، انحراف معیار \lr{11.97}، کمینه‌ی منفیِ
\lr{56.87} و بیشینه‌ی \lr{101.92}. انحراف معیارِ آن به‌روشنی کمتر از انحراف معیارِ سری ساعتی
(\lr{15.94}) است، چون میانگین‌گیری بخشی از نوسان درون‌روزی را حذف می‌کند.

همین کاهش نوسان، پرسش پژوهشیِ چهارم را می‌سازد. اگر هدفِ روزانه ذاتاً هموارتر از
هدف ساعتی است، آیا بهتر است آن را مستقیماً مدل کرد، یا از میانگین‌گیریِ ۲۴ پیش‌بینیِ
ساعتی به دست آورد؟ هر دو مسیر در این پژوهش ساخته و در \cref{sec:4-4} با هم سنجیده می‌شوند.
